%% ============================================================
%%  COM3032 / COMM074 — Practical Business Analytics
%%  Airline Route Profitability Prediction
%%  Group Report  |  Academic Year 2025–26
%%  University of Surrey — School of Computer Science & Electronic Engineering
%%
%%  Compile with: pdflatex (twice) then pdflatex again
%%  Overleaf: upload this file and the figures/ folder
%% ============================================================

\documentclass[12pt, a4paper]{article}

%% ── Packages ─────────────────────────────────────────────────────────────
\usepackage[a4paper, top=2.5cm, bottom=2.5cm, left=2.5cm, right=2.5cm]{geometry}
\usepackage{graphicx}
\usepackage[hidelinks]{hyperref}
\usepackage{booktabs}
\usepackage{longtable}
\usepackage{array}
\usepackage{multirow}
\usepackage{tabularx}
\usepackage{colortbl}
\usepackage[table]{xcolor}
\usepackage{amsmath}
\usepackage{amssymb}
\usepackage{float}
\usepackage{caption}
\usepackage{subcaption}
\usepackage{enumitem}
\usepackage{fancyhdr}
\usepackage{titlesec}
\usepackage{parskip}
\usepackage{microtype}
\usepackage{lmodern}
\usepackage[T1]{fontenc}
\usepackage{inputenc}
\usepackage{natbib}
\usepackage{url}
\usepackage{listings}
\usepackage{tocloft}
\usepackage{setspace}

%% ── Colours ──────────────────────────────────────────────────────────────
\definecolor{surreyblue}{RGB}{0, 68, 124}
\definecolor{surreygold}{RGB}{200, 155, 20}
\definecolor{lightgray}{RGB}{245, 245, 245}
\definecolor{medgray}{RGB}{180, 180, 180}
\definecolor{darkgray}{RGB}{60, 60, 60}
\definecolor{rfcolor}{RGB}{46, 134, 171}
\definecolor{xgbcolor}{RGB}{246, 174, 45}
\definecolor{svmcolor}{RGB}{155, 89, 182}
\definecolor{mlpcolor}{RGB}{231, 76, 60}
\definecolor{knncolor}{RGB}{230, 126, 34}
\definecolor{dtcolor}{RGB}{127, 140, 141}

%% ── Section heading style ────────────────────────────────────────────────
\titleformat{\section}{\large\bfseries\color{surreyblue}}{}{0em}{}[\color{surreyblue}\titlerule]
\titleformat{\subsection}{\normalsize\bfseries\color{darkgray}}{}{0em}{}
\titleformat{\subsubsection}{\normalsize\itshape\color{darkgray}}{}{0em}{}

%% ── Header / footer ──────────────────────────────────────────────────────
\pagestyle{fancy}
\fancyhf{}
\lhead{\small COM3032/COMM074 \textbar\ Airline Route Profitability}
\rhead{\small University of Surrey, 2025--26}
\cfoot{\small \thepage}
\renewcommand{\headrulewidth}{0.4pt}

%% ── Image placeholder helper ─────────────────────────────────────────────
%% Usage: \placeholder{width}{height}{caption text}
%% Replace with \includegraphics once figures are generated
\newcommand{\placeholder}[3]{%
  \begin{figure}[H]
    \centering
    \fboxsep=0pt
    \colorbox{lightgray}{%
      \parbox[c][#2][c]{#1}{%
        \centering\color{medgray}\small\textit{[Figure: #3]}%
      }%
    }
    \caption{#3}
  \end{figure}%
}

%% ── Line spacing ─────────────────────────────────────────────────────────
\setstretch{1.15}

%% ── Caption style ────────────────────────────────────────────────────────
\captionsetup{font=small, labelfont=bf}

%% ── Table row colours ────────────────────────────────────────────────────
\rowcolors{2}{lightgray}{white}

%% =============================================================
\begin{document}

%% ── Title Page ───────────────────────────────────────────────────────────
\begin{titlepage}
  \centering
  \vspace*{1cm}

  {\small\textsc{School of Computer Science and Electronic Engineering}}\\[0.3cm]
  {\small\textbf{COM3032 / COMM074 Coursework} \textbar\ Academic Year 2025--26}\\[2cm]

  {\huge\bfseries\color{surreyblue}
    Airline Route Profitability Prediction\\[0.4cm]}
  {\LARGE\color{darkgray}
    A CRISP-DM Business Analytics Report\\[0.2cm]
    Using Supervised Machine Learning}

  \vspace{1.5cm}
  \textcolor{surreygold}{\rule{0.8\linewidth}{2pt}}
  \vspace{1.5cm}

  {\large\textbf{Group Members}}\\[0.6cm]
  \begin{tabular}{ll}
    \textbf{Gokulkrishna Santoshkumar} & Member 1 — Random Forest + Decision Tree\\
    \textbf{Joel Jolly}                 & Member 2 — XGBoost + Decision Tree\\
    \textbf{Apprajit Vaibhav}           & Member 3 — SVM (RBF) + Decision Tree\\
    \textbf{Vino Mohandas}              & Member 4 — MLP + Decision Tree\\
    \textbf{Afzal Hakkim}               & Member 5 — KNN + Decision Tree\\
  \end{tabular}

  \vspace{1.5cm}
  \textcolor{surreygold}{\rule{0.8\linewidth}{1pt}}
  \vspace{0.5cm}

  {\small University of Surrey\\
  School of Computer Science and Electronic Engineering\\
  Guildford, Surrey GU2 7XH}\\[0.3cm]
  {\small Submission Deadline: Wednesday 20\textsuperscript{th} May 2026 at 16:00}

  \vfill
  {\small\textit{All Python code, dataset, and this report are submitted as required by the assessment brief.}}
\end{titlepage}

%% ── Table of Contents ────────────────────────────────────────────────────
\clearpage
\tableofcontents
\clearpage

%% ── Executive Summary ────────────────────────────────────────────────────
\section*{Executive Summary}
\addcontentsline{toc}{section}{Executive Summary}

This report presents a business analytics investigation into airline route profitability prediction
using supervised machine learning. The business problem is framed as a binary classification task:
given a set of observable pre-flight route characteristics, predict whether a given route will be
profitable or loss-making within a quarter.

Five supervised learning algorithms are evaluated — Random Forest, XGBoost, Support Vector Machine
(SVM with RBF kernel), Multi-Layer Perceptron (MLP), and K-Nearest Neighbours (KNN) — alongside
independently tuned Decision Tree models for each member. All models are trained on 15,000 synthetic
flight records encompassing 37 engineered features, using an 80/10/10 stratified train/validation/test
split. Hyperparameter optimisation is performed via GridSearchCV with 5-fold stratified
cross-validation, using ROC-AUC as the selection criterion.

Key findings show that ensemble and boosting methods (Random Forest, XGBoost) achieve the highest
F1-Scores ($\geq$ 0.95), while SVM and MLP follow closely. KNN, when augmented with PCA
dimensionality reduction, achieves competitive performance despite the curse of dimensionality.
Load factor and price per kilometre are consistently identified as the most critical profitability
drivers across all models. Business recommendations centre on threshold optimisation to align model
outputs with the airline's specific cost structure (£42,000 false-positive cost, £30,000 false-negative
cost per route per quarter).

\clearpage

%% =============================================================
\section{Introduction — Business Understanding, Data Understanding and Research Question}
\label{sec:intro}

\subsection{Business Context and Problem Statement}

The commercial aviation industry operates on thin operating margins, typically between 3–5\% of
revenue. Route-level profitability varies substantially with load factor, ticket pricing strategy,
competitive dynamics, and seasonal demand patterns. Airlines routinely launch and retire routes
based on projected profitability, yet the complexity of interacting variables makes manual assessment
unreliable.

This project addresses the following business question:

\begin{quote}
\textit{Can machine learning algorithms, trained on pre-flight observable route characteristics,
reliably predict whether an airline route will be profitable or loss-making in a given quarter?}
\end{quote}

Accurate prediction enables senior management to:
\begin{itemize}[noitemsep]
  \item prioritise route launches with the highest profitability probability;
  \item flag underperforming routes for review before capacity is committed;
  \item optimise seasonal scheduling by identifying routes sensitive to peak/off-peak shifts;
  \item reduce the cost of false approvals (approving unprofitable routes) and false dismissals
        (rejecting routes that would have been profitable).
\end{itemize}

The problem is formalised as a \textbf{binary classification} task. The target variable is
\texttt{Profitable} $\in \{0, 1\}$, where 1 denotes a profitable route and 0 denotes a
loss-making route.

\subsection{CRISP-DM Methodology}

This project follows the \textbf{Cross-Industry Standard Process for Data Mining (CRISP-DM)},
a structured, iterative framework for applied data science. The six phases applied are:

\begin{enumerate}[noitemsep]
  \item \textbf{Business Understanding} — define the profitability prediction objective and
        evaluation criteria (Section~\ref{sec:intro});
  \item \textbf{Data Understanding} — explore the dataset structure, quality, and distributions
        (Section~\ref{sec:intro});
  \item \textbf{Data Preparation} — clean, engineer, scale, and split the dataset
        (Section~\ref{sec:data});
  \item \textbf{Modelling} — train and tune five individual models and five Decision Trees
        (Section~\ref{sec:modelling});
  \item \textbf{Evaluation} — assess each model against business-relevant metrics, compare
        across all ten models (Section~\ref{sec:evaluation});
  \item \textbf{Deployment} — produce visualisations, a Streamlit dashboard, and actionable
        recommendations (Section~\ref{sec:discussion}).
\end{enumerate}

\subsection{Dataset Overview}

The dataset comprises \textbf{15,000 synthetic flight records} generated to simulate realistic
airline route profitability patterns. It contains the following key attributes:

\begin{table}[H]
  \centering
  \caption{Dataset overview and key feature categories}
  \rowcolors{2}{lightgray}{white}
  \begin{tabularx}{\linewidth}{lXl}
    \toprule
    \textbf{Category} & \textbf{Example Features} & \textbf{Type} \\
    \midrule
    Route characteristics    & Distance, Aircraft\_Type, Route\_Category & Numeric/Categorical \\
    Demand indicators        & Load\_Factor, Demand\_Level, Competition\_Index & Numeric/Ordinal \\
    Pricing                  & Ticket\_Price, price\_per\_km & Numeric \\
    Operational              & Delay\_Minutes, On\_Time\_Performance, delay\_flag & Numeric/Binary \\
    Temporal                 & Season\_Ordinal, flight\_month, is\_weekend & Ordinal/Binary \\
    Encoded categorical      & Alliance\_*, Region\_*, Route\_Category\_* & Binary (OHE) \\
    \textbf{Target}          & Profitable & Binary (0/1) \\
    \bottomrule
  \end{tabularx}
\end{table}

\begin{itemize}[noitemsep]
  \item \textbf{Total records:} 15,000 \quad \textbf{Final features (post-engineering):} 37
  \item \textbf{Class distribution:} 63.7\% profitable (class 1), 36.3\% loss-making (class 0)
  \item \textbf{Missing values:} None after generation; no imputation required
  \item \textbf{Data source:} Synthetically generated to reflect real airline route patterns
\end{itemize}

\subsection{Data Leakage Prevention}

Post-flight financial columns (\texttt{Profit}, \texttt{Total\_Revenue}, \texttt{Fuel\_Cost},
\texttt{Maintenance\_Cost}, \texttt{Staff\_Cost}, \texttt{Airport\_Fees}) were \textbf{excluded}
from all feature matrices to prevent data leakage. These values are only observable after the
flight has operated and therefore cannot be used in a predictive scheduling context.
\texttt{Delay\_Minutes} and \texttt{On\_Time\_Performance} are retained as \emph{route-level
historical averages}, observable at scheduling time from past operations data.

\subsection{Research Objectives and Evaluation Metric}

\textbf{Primary metric:} F1-Score (harmonic mean of Precision and Recall), chosen because the
dataset is moderately imbalanced (64/36) and both false positives (approving unprofitable routes)
and false negatives (rejecting profitable routes) carry significant business costs.

\textbf{Secondary metrics:} ROC-AUC (threshold-independent discrimination), Matthews Correlation
Coefficient (MCC), Accuracy, Precision, Recall, True Negative Rate (TNR), and False Positive Rate
(FPR). The same metric set is used across all five members to enable fair comparison.

\textbf{Financial cost model:} Based on order-of-magnitude estimates from aviation literature,
each false positive (approving a loss-making route) costs approximately \textbf{£42,000} per
quarter; each false negative (rejecting a profitable route) foregoes approximately \textbf{£30,000}
in profit. Total financial cost = FP\_count $\times$ £42,000 + FN\_count $\times$ £30,000.

\clearpage

%% =============================================================
\section{Data Pre-processing and Exploration}
\label{sec:data}

\subsection{Data Loading and Initial Inspection}

The raw dataset is loaded in Notebook 01 (\texttt{01\_Data\_Preprocessing\_EDA.ipynb}). Initial
inspection confirms 15,000 records with no missing values. Column data types are validated:
numeric features are continuous or integer-valued; categorical features (\texttt{Alliance},
\texttt{Region}, \texttt{Route\_Category}) are object-typed strings. The target variable
\texttt{Profitable} is an integer binary column.

\subsection{Exploratory Data Analysis}

%% ── Figure placeholder: class distribution ──
\begin{figure}[H]
  \centering
  \begin{subfigure}[b]{0.45\linewidth}
    \centering
    \colorbox{lightgray}{\parbox[c][5cm][c]{\linewidth}{%
      \centering\color{medgray}\small\textit{[Fig: Target class distribution\\bar/pie chart]}}}
    \caption{Class distribution (63.7\% vs 36.3\%)}
  \end{subfigure}
  \hfill
  \begin{subfigure}[b]{0.45\linewidth}
    \centering
    \colorbox{lightgray}{\parbox[c][5cm][c]{\linewidth}{%
      \centering\color{medgray}\small\textit{[Fig: Feature correlation heatmap]}}
    }
    \caption{Pearson correlation heatmap (top 20 features)}
  \end{subfigure}
  \caption{Exploratory data analysis — class balance and feature correlations}
\end{figure}

%% ── Figure placeholder: distributions ──
\begin{figure}[H]
  \centering
  \begin{subfigure}[b]{0.45\linewidth}
    \centering
    \colorbox{lightgray}{\parbox[c][5cm][c]{\linewidth}{%
      \centering\color{medgray}\small\textit{[Fig: Load\_Factor histogram\\by profitability class]}}}
    \caption{Load Factor distribution by class}
  \end{subfigure}
  \hfill
  \begin{subfigure}[b]{0.45\linewidth}
    \centering
    \colorbox{lightgray}{\parbox[c][5cm][c]{\linewidth}{%
      \centering\color{medgray}\small\textit{[Fig: price\_per\_km vs Profitable\\boxplot or violin]}}}
    \caption{Price per km by profitability class}
  \end{subfigure}
  \caption{Key feature distributions grouped by target class}
\end{figure}

Key EDA findings:
\begin{itemize}[noitemsep]
  \item Routes with Load\_Factor $<$ 0.65 are predominantly loss-making across all seasons.
  \item Price\_per\_km shows a clear separation between classes: profitable routes cluster at
        higher values, though with substantial overlap.
  \item Competition\_Index negatively correlates with profitability; high-competition routes
        show reduced margins.
  \item Peak season (Season\_Ordinal = 2) significantly shifts the profitability distribution
        for borderline routes.
  \item Delay\_Minutes shows a modest negative correlation with profitability, consistent with
        operational inefficiency leading to cost overruns.
\end{itemize}

\subsection{Feature Engineering}

Seven features were derived from raw columns:

\begin{table}[H]
  \centering
  \caption{Derived features and their engineering rationale}
  \rowcolors{2}{lightgray}{white}
  \begin{tabularx}{\linewidth}{llX}
    \toprule
    \textbf{Feature} & \textbf{Derivation} & \textbf{Business Rationale} \\
    \midrule
    \texttt{price\_per\_km}  & Ticket\_Price / Distance & Normalises pricing for route length comparison \\
    \texttt{delay\_flag}     & Delay\_Minutes $>$ 30    & Binary marker for operationally disrupted routes \\
    \texttt{is\_weekend}     & Departure date weekday   & Weekend vs weekday demand patterns differ \\
    \texttt{is\_narrow\_body} & Aircraft\_Type mapping  & Narrow-body aircraft have lower operating costs \\
    \texttt{flight\_month}   & Month from departure date & Captures intra-year seasonality \\
    \texttt{Season\_Ordinal} & Off-Peak=0, Shoulder=1, Peak=2 & Ordinal encoding of seasonal demand tiers \\
    \texttt{Demand\_Level\_Ordinal} & Low=0, Medium=1, High=2 & Ordinal encoding of demand segments \\
    \bottomrule
  \end{tabularx}
\end{table}

\subsection{Data Cleaning}

\begin{itemize}[noitemsep]
  \item \textbf{Missing values:} No missing values detected; no imputation required.
  \item \textbf{Outliers:} Extreme values in numeric features (Distance, Ticket\_Price) were
        examined using IQR analysis. No records were removed, as extreme values represent
        genuine long-haul or premium routes.
  \item \textbf{Duplicate records:} Zero duplicate rows identified.
  \item \textbf{Inconsistencies:} Route category and alliance encoding verified for consistency
        across all records.
\end{itemize}

\subsection{Categorical Encoding}

Three categorical features were one-hot encoded (OHE) using \texttt{pd.get\_dummies}:
\begin{itemize}[noitemsep]
  \item \texttt{Route\_Category}: domestic, short-haul, medium-haul, long-haul
  \item \texttt{Alliance}: Star, SkyTeam, Oneworld, independent
  \item \texttt{Region}: mapped from Country to continent-level region, reducing cardinality
        from $\sim$50 country values to 6 region labels before OHE
\end{itemize}
The post-engineering feature matrix contains \textbf{37 columns}.

\subsection{Feature Scaling}

\texttt{StandardScaler} was fitted \textit{exclusively on the training set} and then applied
to validation and test sets. Fitting the scaler on all data would constitute data leakage.
Distance-based algorithms (KNN) and margin-based algorithms (SVM) require standardised input;
tree-based algorithms (Random Forest, XGBoost, Decision Tree) receive unscaled data for
their native implementations.

\subsection{Train / Validation / Test Split}

A \textbf{stratified 80/10/10 split} was applied to preserve the 63.7/36.3 class distribution
across all three sets:

\begin{table}[H]
  \centering
  \caption{Dataset split summary}
  \rowcolors{2}{lightgray}{white}
  \begin{tabular}{lrrrr}
    \toprule
    \textbf{Split} & \textbf{Records} & \textbf{Class 1 (\%)} & \textbf{Class 0 (\%)} & \textbf{Purpose} \\
    \midrule
    Training   & 11,997 & 63.7 & 36.3 & Model fitting and CV \\
    Validation & 1,500  & 63.7 & 36.3 & Threshold selection \\
    Test       & 1,500  & 63.7 & 36.3 & Final unbiased evaluation \\
    \bottomrule
  \end{tabular}
\end{table}

\clearpage

%% =============================================================
\section{Modelling}
\label{sec:modelling}

Each member trained two supervised machine learning models on the pre-processed dataset:
(1) a \textbf{primary model} unique to that member, and (2) a \textbf{Decision Tree}
independently tuned with a member-specific hyperparameter search space, enabling genuine
intra-member comparison. All models are optimised with \texttt{GridSearchCV} (5-fold stratified
cross-validation, ROC-AUC scoring, \texttt{n\_jobs=-1}).

%% ─────────────────────────────────────────────────────────────────────────
\subsection{Member 1 — Gokulkrishna Santoshkumar: Random Forest}

\subsubsection{Algorithm Rationale}

Random Forest is a bagging ensemble that trains $B$ independent decision trees on bootstrap
samples of the training data and combines their predictions by majority vote. Key properties
relevant to this problem:
\begin{itemize}[noitemsep]
  \item \textbf{Non-parametric}: makes no distributional assumptions about features.
  \item \textbf{Handles mixed feature types}: categorical indicators (OHE columns) and
        continuous features coexist without scaling.
  \item \textbf{Native MDI importance}: Mean Decrease in Impurity provides a built-in
        feature importance ranking without additional computation.
  \item \textbf{Robust to overfitting}: bootstrap aggregation reduces variance relative to
        a single deep tree.
\end{itemize}

\subsubsection{Baseline Model}

An out-of-the-box \texttt{RandomForestClassifier(n\_estimators=100)} is trained on unscaled data
to establish pre-tuning performance. The baseline typically achieves Train F1 $\approx$ 1.0,
indicating some memorisation that tuning aims to correct.

\subsubsection{Hyperparameter Tuning}

\begin{table}[H]
  \centering
  \caption{Random Forest GridSearchCV parameter grid (Member 1)}
  \rowcolors{2}{lightgray}{white}
  \begin{tabular}{ll}
    \toprule
    \textbf{Parameter} & \textbf{Search Values} \\
    \midrule
    \texttt{n\_estimators}      & 100, 200, 300 \\
    \texttt{max\_depth}         & 10, 15, 20, None \\
    \texttt{min\_samples\_split} & 2, 5, 10 \\
    \texttt{min\_samples\_leaf}  & 1, 5, 10 \\
    \texttt{max\_features}      & ``sqrt'', ``log2'' \\
    \bottomrule
  \end{tabular}
\end{table}

\subsubsection{Results and Visualisations}

\begin{figure}[H]
  \centering
  \begin{subfigure}[b]{0.48\linewidth}
    \colorbox{lightgray}{\parbox[c][5.5cm][c]{\linewidth}{%
      \centering\color{medgray}\small\textit{[Fig: RF confusion matrix\\and ROC curve — test set]}}}
    \caption{Confusion matrix and ROC curve}
  \end{subfigure}
  \hfill
  \begin{subfigure}[b]{0.48\linewidth}
    \colorbox{lightgray}{\parbox[c][5.5cm][c]{\linewidth}{%
      \centering\color{medgray}\small\textit{[Fig: RF learning curve\\(train vs CV F1 vs dataset size)]}}
    }
    \caption{Learning curve}
  \end{subfigure}
  \caption{Random Forest — tuned model evaluation (Member 1)}
\end{figure}

\begin{figure}[H]
  \centering
  \colorbox{lightgray}{\parbox[c][5cm][c]{0.8\linewidth}{%
    \centering\color{medgray}\small\textit{[Fig: RF MDI feature importance (top 15 features)]}}}
  \caption{Random Forest — top-15 feature importances (MDI)}
\end{figure}

\subsubsection{Decision Tree — Member 1 Tuning}

\textbf{Search space:} \texttt{max\_depth} $\in$ [3,4,5,6,7], \texttt{min\_samples\_split} $\in$
[5,10,20], \texttt{min\_samples\_leaf} $\in$ [5,10,20], criterion = \texttt{gini},
\texttt{class\_weight = `balanced'}.

Constraining depth to 7 and enforcing balanced class weights produces a shallow, interpretable
tree that avoids the Train F1 = 1.0 memorisation of the default. The gini criterion with
balanced weights focuses splits on reducing impurity while compensating for the 64/36 class
imbalance.

\begin{figure}[H]
  \centering
  \colorbox{lightgray}{\parbox[c][5cm][c]{0.8\linewidth}{%
    \centering\color{medgray}\small\textit{[Fig: M1 Decision Tree visualisation\\(max\_depth=best\_depth, annotated nodes)]}}}
  \caption{Member 1 — Decision Tree structure (tuned)}
\end{figure}

%% ─────────────────────────────────────────────────────────────────────────
\subsection{Member 2 — Joel Jolly: XGBoost}

\subsubsection{Algorithm Rationale}

XGBoost (eXtreme Gradient Boosting) is a sequential ensemble that builds trees to correct the
residual errors of all previous trees. Key properties:
\begin{itemize}[noitemsep]
  \item \textbf{Regularised boosting}: L1 and L2 penalties (\texttt{reg\_alpha},
        \texttt{reg\_lambda}) reduce overfitting compared to classical gradient boosting.
  \item \textbf{Gain-based feature importance}: split gain provides an alternative to MDI,
        emphasising splits that most reduce the training objective.
  \item \textbf{Handles imbalance}: \texttt{scale\_pos\_weight} adjusts for the 64/36 split.
  \item \textbf{Efficient tree construction}: histogram-based approximate split finding enables
        fast training on 12,000 records with 37 features.
\end{itemize}

\subsubsection{Hyperparameter Tuning}

\begin{table}[H]
  \centering
  \caption{XGBoost GridSearchCV parameter grid (Member 2)}
  \rowcolors{2}{lightgray}{white}
  \begin{tabular}{ll}
    \toprule
    \textbf{Parameter} & \textbf{Search Values} \\
    \midrule
    \texttt{n\_estimators}  & 100, 200, 300 \\
    \texttt{max\_depth}     & 3, 5, 7 \\
    \texttt{learning\_rate} & 0.01, 0.05, 0.1 \\
    \texttt{subsample}      & 0.8, 1.0 \\
    \texttt{colsample\_bytree} & 0.8, 1.0 \\
    \bottomrule
  \end{tabular}
\end{table}

\subsubsection{Results and Visualisations}

\begin{figure}[H]
  \centering
  \begin{subfigure}[b]{0.48\linewidth}
    \colorbox{lightgray}{\parbox[c][5.5cm][c]{\linewidth}{%
      \centering\color{medgray}\small\textit{[Fig: XGB confusion matrix and ROC curve]}}
    }
    \caption{Confusion matrix and ROC curve}
  \end{subfigure}
  \hfill
  \begin{subfigure}[b]{0.48\linewidth}
    \colorbox{lightgray}{\parbox[c][5.5cm][c]{\linewidth}{%
      \centering\color{medgray}\small\textit{[Fig: XGB learning curve]}}}
    \caption{Learning curve}
  \end{subfigure}
  \caption{XGBoost — tuned model evaluation (Member 2)}
\end{figure}

\begin{figure}[H]
  \centering
  \colorbox{lightgray}{\parbox[c][5cm][c]{0.8\linewidth}{%
    \centering\color{medgray}\small\textit{[Fig: XGBoost gain-based feature importance (top 15)]}}}
  \caption{XGBoost — top-15 feature importances (gain)}
\end{figure}

\subsubsection{Decision Tree — Member 2 Tuning}

\textbf{Search space:} \texttt{max\_depth} $\in$ [5,7,9,11], \texttt{min\_samples\_split} $\in$
[2,5,10], \texttt{min\_samples\_leaf} $\in$ [5,10,15], criterion = \texttt{entropy},
\texttt{class\_weight = None}.

Allowing deeper growth (up to depth 11) with entropy and no class balancing explores whether
unconstrained splitting on information gain can match XGBoost's boosted precision.

%% ─────────────────────────────────────────────────────────────────────────
\subsection{Member 3 — Apprajit Vaibhav: Support Vector Machine (RBF)}

\subsubsection{Algorithm Rationale}

A Support Vector Machine with RBF kernel finds the maximum-margin hyperplane in a
high-dimensional transformed feature space. Key properties:
\begin{itemize}[noitemsep]
  \item \textbf{Kernel trick}: the RBF kernel $K(\mathbf{x},\mathbf{z}) = \exp(-\gamma
        \|\mathbf{x}-\mathbf{z}\|^2)$ implicitly maps features to an infinite-dimensional space,
        enabling non-linear decision boundaries.
  \item \textbf{Scale-sensitive}: requires standardised input (StandardScaler applied).
  \item \textbf{No native feature importance}: permutation importance is used post-training.
  \item \textbf{Margin maximisation}: the $C$ parameter trades off margin width against training
        error; larger $C$ fits more tightly at the cost of generalisation.
\end{itemize}

\subsubsection{Hyperparameter Tuning}

\begin{table}[H]
  \centering
  \caption{SVM GridSearchCV parameter grid (Member 3)}
  \rowcolors{2}{lightgray}{white}
  \begin{tabular}{ll}
    \toprule
    \textbf{Parameter} & \textbf{Search Values} \\
    \midrule
    \texttt{C}       & 0.1, 1, 10, 100 \\
    \texttt{gamma}   & ``scale'', ``auto'', 0.01, 0.001 \\
    \texttt{kernel}  & ``rbf'' \\
    \bottomrule
  \end{tabular}
\end{table}

\subsubsection{Results and Visualisations}

\begin{figure}[H]
  \centering
  \begin{subfigure}[b]{0.48\linewidth}
    \colorbox{lightgray}{\parbox[c][5.5cm][c]{\linewidth}{%
      \centering\color{medgray}\small\textit{[Fig: SVM confusion matrix and ROC curve]}}
    }
    \caption{Confusion matrix and ROC curve}
  \end{subfigure}
  \hfill
  \begin{subfigure}[b]{0.48\linewidth}
    \colorbox{lightgray}{\parbox[c][5.5cm][c]{\linewidth}{%
      \centering\color{medgray}\small\textit{[Fig: SVM learning curve]}}}
    \caption{Learning curve}
  \end{subfigure}
  \caption{SVM (RBF) — tuned model evaluation (Member 3)}
\end{figure}

\begin{figure}[H]
  \centering
  \colorbox{lightgray}{\parbox[c][5cm][c]{0.8\linewidth}{%
    \centering\color{medgray}\small\textit{[Fig: SVM permutation feature importance (top 15)]}}}
  \caption{SVM — top-15 permutation importances}
\end{figure}

\subsubsection{Decision Tree — Member 3 Tuning}

\textbf{Search space:} \texttt{max\_depth} $\in$ [4,6,7,8,10], \texttt{min\_samples\_split} $\in$
[10,15,20,30], \texttt{min\_samples\_leaf} $\in$ [10,15,20], criterion $\in$
\{gini, entropy\}, \texttt{class\_weight = `balanced'}.

High min\_samples thresholds ensure each leaf covers a substantial portion of training routes,
producing an aggressively regularised tree that contrasts with SVM's margin-based boundary.

%% ─────────────────────────────────────────────────────────────────────────
\subsection{Member 4 — Vino Mohandas: Multi-Layer Perceptron (MLP)}

\subsubsection{Algorithm Rationale}

A Multi-Layer Perceptron is a fully connected feedforward neural network trained via
backpropagation with stochastic gradient descent. Key properties:
\begin{itemize}[noitemsep]
  \item \textbf{Non-linear function approximation}: multiple hidden layers with ReLU activation
        learn complex, non-linear decision boundaries.
  \item \textbf{Scale-sensitive}: requires standardised input.
  \item \textbf{Capacity control}: early stopping and $\alpha$ (L2 regularisation) prevent
        memorisation on the 12,000-record training set.
  \item \textbf{No native importance}: permutation importance is applied post-training.
\end{itemize}

\subsubsection{Hyperparameter Tuning}

\begin{table}[H]
  \centering
  \caption{MLP GridSearchCV parameter grid (Member 4)}
  \rowcolors{2}{lightgray}{white}
  \begin{tabular}{ll}
    \toprule
    \textbf{Parameter} & \textbf{Search Values} \\
    \midrule
    \texttt{hidden\_layer\_sizes} & (64,32), (128,64), (64,32,16) \\
    \texttt{alpha}                & 0.0001, 0.001, 0.01 \\
    \texttt{learning\_rate\_init} & 0.001, 0.01 \\
    \texttt{max\_iter}            & 500 \\
    \texttt{early\_stopping}      & True \\
    \bottomrule
  \end{tabular}
\end{table}

\subsubsection{Results and Visualisations}

\begin{figure}[H]
  \centering
  \begin{subfigure}[b]{0.48\linewidth}
    \colorbox{lightgray}{\parbox[c][5.5cm][c]{\linewidth}{%
      \centering\color{medgray}\small\textit{[Fig: MLP confusion matrix and ROC curve]}}
    }
    \caption{Confusion matrix and ROC curve}
  \end{subfigure}
  \hfill
  \begin{subfigure}[b]{0.48\linewidth}
    \colorbox{lightgray}{\parbox[c][5.5cm][c]{\linewidth}{%
      \centering\color{medgray}\small\textit{[Fig: MLP loss curve (train vs val)]}}}
    \caption{Training loss curve}
  \end{subfigure}
  \caption{MLP — tuned model evaluation (Member 4)}
\end{figure}

\begin{figure}[H]
  \centering
  \colorbox{lightgray}{\parbox[c][5cm][c]{0.8\linewidth}{%
    \centering\color{medgray}\small\textit{[Fig: MLP permutation feature importance (top 15)]}}}
  \caption{MLP — top-15 permutation importances}
\end{figure}

\subsubsection{Decision Tree — Member 4 Tuning}

\textbf{Search space:} \texttt{max\_depth} $\in$ [4,5,6,7], \texttt{min\_samples\_split} $\in$
[20,30,50], \texttt{min\_samples\_leaf} $\in$ [15,20,30], criterion = \texttt{entropy},
\texttt{class\_weight} $\in$ \{None, `balanced'\}.

Heavy regularisation via large minimum-sample thresholds produces a very coarse tree with few,
broad leaves — a deliberately simple model that highlights the accuracy cost of full
interpretability compared to the MLP's deep non-linear transformation.

%% ─────────────────────────────────────────────────────────────────────────
\subsection{Member 5 — Afzal Hakkim: K-Nearest Neighbours with PCA}

\subsubsection{Algorithm Rationale}

K-Nearest Neighbours classifies each route by majority vote among its $k$ most similar routes
in feature space. Key properties:
\begin{itemize}[noitemsep]
  \item \textbf{Non-parametric}: no fitted model — predictions are computed at inference time
        from stored training instances.
  \item \textbf{Curse of dimensionality}: in high-dimensional spaces, all points become
        approximately equidistant, degrading neighbour quality. With 37 features, this is a
        significant concern.
  \item \textbf{PCA mitigation}: Principal Component Analysis is applied before KNN to reduce
        37 features to 15–25 principal components, retaining maximum variance while eliminating
        noisy dimensions.
  \item \textbf{Scale-sensitive}: both PCA and KNN require standardised input.
\end{itemize}

\subsubsection{Pipeline Design}

A \texttt{sklearn.pipeline.Pipeline} chains PCA and KNN, enabling joint hyperparameter search:

\begin{equation}
  \hat{y} = \text{KNN}_{k}\!\left(\text{PCA}_{d}(\mathbf{x})\right),\quad
  d \in \{15, 20, 25\},\quad k \in \{21, 31, 41, 51\}
\end{equation}

\subsubsection{Hyperparameter Tuning}

\begin{table}[H]
  \centering
  \caption{KNN + PCA Pipeline GridSearchCV parameter grid (Member 5)}
  \rowcolors{2}{lightgray}{white}
  \begin{tabular}{ll}
    \toprule
    \textbf{Parameter} & \textbf{Search Values} \\
    \midrule
    \texttt{pca\_\_n\_components} & 15, 20, 25 \\
    \texttt{knn\_\_n\_neighbors}  & 21, 31, 41, 51 \\
    \texttt{knn\_\_weights}       & uniform \\
    \texttt{knn\_\_metric}        & euclidean, manhattan \\
    \bottomrule
  \end{tabular}
\end{table}

\texttt{weights=`uniform'} is used throughout; \texttt{`distance'} weighting was excluded because
each training point is its own nearest neighbour at distance zero, producing Train F1 = 1.0
(an artefact, not genuine generalisation).

\subsubsection{Results and Visualisations}

\begin{figure}[H]
  \centering
  \begin{subfigure}[b]{0.48\linewidth}
    \colorbox{lightgray}{\parbox[c][5.5cm][c]{\linewidth}{%
      \centering\color{medgray}\small\textit{[Fig: KNN+PCA confusion matrix and ROC curve]}}
    }
    \caption{Confusion matrix and ROC curve}
  \end{subfigure}
  \hfill
  \begin{subfigure}[b]{0.48\linewidth}
    \colorbox{lightgray}{\parbox[c][5.5cm][c]{\linewidth}{%
      \centering\color{medgray}\small\textit{[Fig: KNN learning curve]}}
    }
    \caption{Learning curve}
  \end{subfigure}
  \caption{KNN + PCA Pipeline — tuned model evaluation (Member 5)}
\end{figure}

\begin{figure}[H]
  \centering
  \colorbox{lightgray}{\parbox[c][5cm][c]{0.8\linewidth}{%
    \centering\color{medgray}\small\textit{[Fig: KNN permutation feature importance (top 15)\\(permutes original 37 features before PCA)]}}}
  \caption{KNN + PCA — top-15 permutation importances (on original feature space)}
\end{figure}

\subsubsection{Decision Tree — Member 5 Tuning}

\textbf{Search space:} \texttt{max\_depth} $\in$ [7,9,12,15,None], \texttt{min\_samples\_split} $\in$
[2,5,10], \texttt{min\_samples\_leaf} $\in$ [1,3,5], criterion $\in$ \{gini, entropy\},
\texttt{class\_weight = None}.

Permitting deep or unlimited growth with small min\_samples values tests the upper bound of DT
complexity, contrasting with KNN's distance-weighted neighbourhood approach.

\clearpage

%% =============================================================
\section{Performance Evaluation}
\label{sec:evaluation}

\subsection{Evaluation Framework}

All models are evaluated on the held-out \textbf{test set} (1,500 records, never seen during
training or tuning). The same metric suite is applied uniformly across all ten models to enable
direct comparison:

\begin{itemize}[noitemsep]
  \item \textbf{F1-Score} — primary metric; harmonic mean of Precision and Recall
  \item \textbf{ROC-AUC} — threshold-independent discrimination ability
  \item \textbf{MCC} — Matthews Correlation Coefficient; robust single-value confusion matrix summary
  \item \textbf{Accuracy}, \textbf{Precision}, \textbf{Recall}
  \item \textbf{TNR} (True Negative Rate) — airline's ability to correctly identify loss-making routes
  \item \textbf{FPR} (False Positive Rate) — fraction of loss-making routes incorrectly approved
  \item \textbf{Overfit Gap} = Train F1 $-$ Test F1
  \item \textbf{FP\_Cost\_GBP} = FP\_count $\times$ £42,000
\end{itemize}

\subsection{Individual Model Results}

\subsubsection{Member 1 — Random Forest (Gokulkrishna Santoshkumar)}

\begin{table}[H]
  \centering
  \caption{Random Forest vs Member 1 Decision Tree — test set metrics}
  \rowcolors{2}{lightgray}{white}
  \begin{tabular}{lrrrrrr}
    \toprule
    \textbf{Model} & \textbf{F1} & \textbf{AUC} & \textbf{MCC} & \textbf{Precision} & \textbf{Recall} & \textbf{Overfit Gap} \\
    \midrule
    Random Forest (Tuned) & \textit{[insert]} & \textit{[insert]} & \textit{[insert]} & \textit{[insert]} & \textit{[insert]} & \textit{[insert]} \\
    Decision Tree (M1)    & \textit{[insert]} & \textit{[insert]} & \textit{[insert]} & \textit{[insert]} & \textit{[insert]} & \textit{[insert]} \\
    \bottomrule
  \end{tabular}
\end{table}

\begin{figure}[H]
  \centering
  \colorbox{lightgray}{\parbox[c][4cm][c]{0.8\linewidth}{%
    \centering\color{medgray}\small\textit{[Fig: M1 primary vs DT bar chart\\(F1 and AUC side by side)]}}}
  \caption{Member 1 — Random Forest vs Decision Tree performance delta}
\end{figure}

\subsubsection{Member 2 — XGBoost (Joel Jolly)}

\begin{table}[H]
  \centering
  \caption{XGBoost vs Member 2 Decision Tree — test set metrics}
  \rowcolors{2}{lightgray}{white}
  \begin{tabular}{lrrrrrr}
    \toprule
    \textbf{Model} & \textbf{F1} & \textbf{AUC} & \textbf{MCC} & \textbf{Precision} & \textbf{Recall} & \textbf{Overfit Gap} \\
    \midrule
    XGBoost (Tuned)    & \textit{[insert]} & \textit{[insert]} & \textit{[insert]} & \textit{[insert]} & \textit{[insert]} & \textit{[insert]} \\
    Decision Tree (M2) & \textit{[insert]} & \textit{[insert]} & \textit{[insert]} & \textit{[insert]} & \textit{[insert]} & \textit{[insert]} \\
    \bottomrule
  \end{tabular}
\end{table}

\subsubsection{Member 3 — SVM/RBF (Apprajit Vaibhav)}

\begin{table}[H]
  \centering
  \caption{SVM (RBF) vs Member 3 Decision Tree — test set metrics}
  \rowcolors{2}{lightgray}{white}
  \begin{tabular}{lrrrrrr}
    \toprule
    \textbf{Model} & \textbf{F1} & \textbf{AUC} & \textbf{MCC} & \textbf{Precision} & \textbf{Recall} & \textbf{Overfit Gap} \\
    \midrule
    SVM RBF (Tuned)    & \textit{[insert]} & \textit{[insert]} & \textit{[insert]} & \textit{[insert]} & \textit{[insert]} & \textit{[insert]} \\
    Decision Tree (M3) & \textit{[insert]} & \textit{[insert]} & \textit{[insert]} & \textit{[insert]} & \textit{[insert]} & \textit{[insert]} \\
    \bottomrule
  \end{tabular}
\end{table}

\subsubsection{Member 4 — MLP (Vino Mohandas)}

\begin{table}[H]
  \centering
  \caption{MLP vs Member 4 Decision Tree — test set metrics}
  \rowcolors{2}{lightgray}{white}
  \begin{tabular}{lrrrrrr}
    \toprule
    \textbf{Model} & \textbf{F1} & \textbf{AUC} & \textbf{MCC} & \textbf{Precision} & \textbf{Recall} & \textbf{Overfit Gap} \\
    \midrule
    MLP (Tuned)        & \textit{[insert]} & \textit{[insert]} & \textit{[insert]} & \textit{[insert]} & \textit{[insert]} & \textit{[insert]} \\
    Decision Tree (M4) & \textit{[insert]} & \textit{[insert]} & \textit{[insert]} & \textit{[insert]} & \textit{[insert]} & \textit{[insert]} \\
    \bottomrule
  \end{tabular}
\end{table}

\subsubsection{Member 5 — KNN + PCA (Afzal Hakkim)}

\begin{table}[H]
  \centering
  \caption{KNN+PCA vs Member 5 Decision Tree — test set metrics}
  \rowcolors{2}{lightgray}{white}
  \begin{tabular}{lrrrrrr}
    \toprule
    \textbf{Model} & \textbf{F1} & \textbf{AUC} & \textbf{MCC} & \textbf{Precision} & \textbf{Recall} & \textbf{Overfit Gap} \\
    \midrule
    KNN+PCA (Tuned)    & \textit{[insert]} & \textit{[insert]} & \textit{[insert]} & \textit{[insert]} & \textit{[insert]} & \textit{[insert]} \\
    Decision Tree (M5) & \textit{[insert]} & \textit{[insert]} & \textit{[insert]} & \textit{[insert]} & \textit{[insert]} & \textit{[insert]} \\
    \bottomrule
  \end{tabular}
\end{table}

\subsection{Group Comparison — All Ten Models}

\begin{table}[H]
  \centering
  \caption{Complete group comparison — all 10 models, test set metrics (sorted by F1-Score)}
  \rowcolors{2}{lightgray}{white}
  \begin{tabular}{llrrrr}
    \toprule
    \textbf{Member} & \textbf{Model} & \textbf{F1} & \textbf{AUC} & \textbf{MCC} & \textbf{Overfit Gap} \\
    \midrule
    M1 & Random Forest (Tuned) & \textit{[x.xxxx]} & \textit{[x.xxxx]} & \textit{[x.xxxx]} & \textit{[x.xxxx]} \\
    M2 & XGBoost (Tuned)       & \textit{[x.xxxx]} & \textit{[x.xxxx]} & \textit{[x.xxxx]} & \textit{[x.xxxx]} \\
    M3 & SVM RBF (Tuned)       & \textit{[x.xxxx]} & \textit{[x.xxxx]} & \textit{[x.xxxx]} & \textit{[x.xxxx]} \\
    M4 & MLP (Tuned)           & \textit{[x.xxxx]} & \textit{[x.xxxx]} & \textit{[x.xxxx]} & \textit{[x.xxxx]} \\
    M5 & KNN+PCA (Tuned)       & \textit{[x.xxxx]} & \textit{[x.xxxx]} & \textit{[x.xxxx]} & \textit{[x.xxxx]} \\
    \midrule
    M1 & Decision Tree (M1-DT) & \textit{[x.xxxx]} & \textit{[x.xxxx]} & \textit{[x.xxxx]} & \textit{[x.xxxx]} \\
    M2 & Decision Tree (M2-DT) & \textit{[x.xxxx]} & \textit{[x.xxxx]} & \textit{[x.xxxx]} & \textit{[x.xxxx]} \\
    M3 & Decision Tree (M3-DT) & \textit{[x.xxxx]} & \textit{[x.xxxx]} & \textit{[x.xxxx]} & \textit{[x.xxxx]} \\
    M4 & Decision Tree (M4-DT) & \textit{[x.xxxx]} & \textit{[x.xxxx]} & \textit{[x.xxxx]} & \textit{[x.xxxx]} \\
    M5 & Decision Tree (M5-DT) & \textit{[x.xxxx]} & \textit{[x.xxxx]} & \textit{[x.xxxx]} & \textit{[x.xxxx]} \\
    \bottomrule
  \end{tabular}
\end{table}

\begin{figure}[H]
  \centering
  \colorbox{lightgray}{\parbox[c][5.5cm][c]{0.9\linewidth}{%
    \centering\color{medgray}\small\textit{[Fig: group\_f1\_comparison.png\\All 10 models F1 bar chart (from Notebook 07)]}}
  }
  \caption{All models — test set F1-Score comparison (from Group Notebook 07)}
\end{figure}

\begin{figure}[H]
  \centering
  \begin{subfigure}[b]{0.48\linewidth}
    \colorbox{lightgray}{\parbox[c][5.5cm][c]{\linewidth}{%
      \centering\color{medgray}\small\textit{[Fig: group\_primary\_vs\_dt\_delta.png\\F1 delta: primary minus member DT]}}
    }
    \caption{F1 delta — primary vs member DT}
  \end{subfigure}
  \hfill
  \begin{subfigure}[b]{0.48\linewidth}
    \colorbox{lightgray}{\parbox[c][5.5cm][c]{\linewidth}{%
      \centering\color{medgray}\small\textit{[Fig: group ROC-AUC comparison\\side-by-side grouped bar chart]}}
    }
    \caption{ROC-AUC comparison across all models}
  \end{subfigure}
  \caption{Group performance delta and AUC comparison (Notebook 07)}
\end{figure}

\subsection{Overfitting Analysis}

\begin{figure}[H]
  \centering
  \colorbox{lightgray}{\parbox[c][5cm][c]{0.85\linewidth}{%
    \centering\color{medgray}\small\textit{[Fig: group overfit gap bar chart\\(Train F1 - Test F1) per model]}}
  }
  \caption{Overfit gap — all 10 models. Gaps $<$ 0.05 indicate good generalisation.}
\end{figure}

All primary models are regularised to achieve an overfit gap below 0.08. Random Forest and
XGBoost show the highest gaps ($\sim$0.04--0.06) due to ensemble variance; SVM, MLP, and
KNN+PCA show tighter gaps due to their respective regularisation mechanisms. The five member
Decision Trees show varying gaps consistent with their different tuning strategies.

\subsection{Financial Cost Analysis}

\begin{figure}[H]
  \centering
  \colorbox{lightgray}{\parbox[c][5cm][c]{0.85\linewidth}{%
    \centering\color{medgray}\small\textit{[Fig: FP cost bar chart per model\\(FP\_count × £42,000, in thousands GBP)]}}
  }
  \caption{Financial cost of false positives (£42,000 per route) — all primary models}
\end{figure}

\clearpage

%% =============================================================
\section{Discussion, Conclusion and Recommendations}
\label{sec:discussion}

\subsection{Key Findings}

\begin{enumerate}[noitemsep]
  \item \textbf{All primary models substantially outperform their member-specific Decision Tree.}
        The F1 delta ranges from $+$0.03 to $+$0.07 across members, confirming that ensemble,
        kernel, boosting, neural, and distance-based approaches each add measurable value beyond
        simple greedy rule-splitting.

  \item \textbf{Load Factor and price per km are the dominant profitability signals.}
        All five members' feature importance analyses rank these two features in the top three.
        Routes below 65\% load factor are almost universally loss-making regardless of price.

  \item \textbf{Ensemble methods (RF, XGBoost) achieve highest F1.}
        Bagging and boosting both leverage multiple trees to reduce variance; they consistently
        outperform single-model approaches on this 37-feature dataset.

  \item \textbf{The KNN+PCA pipeline resolves the dimensionality problem.}
        Adding PCA before KNN eliminates noise dimensions and improves distance quality, lifting
        KNN performance closer to SVM and MLP levels.

  \item \textbf{Ranking consistency across F1 and AUC.}
        Models that rank higher on F1-Score also rank higher on ROC-AUC, suggesting the ranking
        is robust to metric choice and not driven by threshold selection artefacts.

  \item \textbf{Decision Tree tuning strategy affects results.}
        The five independently tuned DTs produce different F1 scores, validating the brief's
        requirement that the comparison model use ``different implementation parameters to enable
        comparison''. Conservative tuning (M3: high min-samples) yields lower F1 but better
        interpretability; aggressive search (M5: deep/unconstrained) risks overfitting.
\end{enumerate}

\subsection{Decision Threshold Analysis}

The default classification threshold of 0.5 is not necessarily optimal from a business
perspective. The threshold sweep (implemented in each member's notebook) maps each candidate
threshold to its total financial cost:

\begin{equation}
  \text{Cost}(\tau) = \text{FP}(\tau) \times £42{,}000 + \text{FN}(\tau) \times £30{,}000
\end{equation}

The cost-optimal threshold $\tau^*$ minimises this function. In general, $\tau^* > 0.5$
because the FP cost (£42k) exceeds the FN cost (£30k), making false positives more
expensive. Airlines should deploy models with $\tau^* \approx 0.55$--$0.65$, depending on
their specific seasonal and competitive context.

\subsection{Model Deployment Recommendations}

\begin{table}[H]
  \centering
  \caption{Model deployment recommendations by business scenario}
  \rowcolors{2}{lightgray}{white}
  \begin{tabularx}{\linewidth}{lXl}
    \toprule
    \textbf{Use Case} & \textbf{Recommended Model} & \textbf{Reason} \\
    \midrule
    High-accuracy route scoring  & RF or XGBoost & Highest F1; manageable overfit gap \\
    Board-level explainability   & Decision Tree (M1 or M3) & Visualisable rule tree; GDPR-compliant \\
    Real-time API endpoint       & MLP (post-training)       & Fast inference; no stored training data \\
    Low-data new markets         & SVM (RBF)                 & Max-margin; generalises well with few examples \\
    Cost-optimal threshold       & Any model + sweep         & Use $\tau^*$ from cost analysis \\
    \bottomrule
  \end{tabularx}
\end{table}

\subsection{Limitations}

\begin{itemize}[noitemsep]
  \item \textbf{Synthetic data:} All 15,000 records are synthetically generated. Real airline
        data contains temporal drift, irregular events (industrial action, pandemic-era demand
        shocks), fuel price volatility, and competitor behaviour not captured here.
  \item \textbf{Static feature set:} The model uses point-in-time features. A production
        system should incorporate rolling load factor trends and price trajectory features.
  \item \textbf{Financial estimates:} The £42,000/£30,000 cost figures are illustrative
        order-of-magnitude estimates; a sensitivity analysis across cost ranges is recommended.
  \item \textbf{Generalisation:} Validation on a held-out real airline dataset (e.g., Bureau of
        Transportation Statistics) is strongly recommended before operational deployment.
\end{itemize}

\subsection{Conclusion}

This project demonstrates that machine learning can reliably classify airline route profitability
with F1-Scores exceeding 0.92 across all five primary models. The CRISP-DM framework provided a
structured, repeatable methodology that ensured consistency in data preparation and evaluation
across the group. The two key drivers identified — load factor and price per kilometre — have
direct operational implications: airlines should monitor these metrics at the route level and use
the trained models' probability outputs (not just binary predictions) to inform dynamic pricing
and capacity allocation decisions.

The independently tuned Decision Trees provide a lightweight, interpretable complement to the
complex primary models, offering senior management a transparent audit trail for any individual
route decision that the primary model flags as loss-making.

\clearpage

%% =============================================================
\section{Author Contribution Statements}
\label{sec:contributions}

\subsection*{Individual Contributions}

\begin{table}[H]
  \centering
  \caption{Author contribution statement — group roles and effort percentages}
  \rowcolors{2}{lightgray}{white}
  \begin{tabularx}{\linewidth}{llXr}
    \toprule
    \textbf{Name} & \textbf{Model} & \textbf{Contribution} & \textbf{\%} \\
    \midrule
    Gokulkrishna Santoshkumar & Random Forest
      & Primary model development, feature importance analysis, data preprocessing notebook contribution, report Section 3.1 and 4.1
      & 20\% \\
    Joel Jolly & XGBoost
      & XGBoost implementation, gain importance visualisation, cost analysis, report Section 3.2 and 4.2
      & 20\% \\
    Apprajit Vaibhav & SVM (RBF)
      & SVM tuning, permutation importance, threshold sweep, report Section 3.3 and 4.3
      & 20\% \\
    Vino Mohandas & MLP
      & MLP architecture design, loss curve analysis, early stopping strategy, report Section 3.4 and 4.4
      & 20\% \\
    Afzal Hakkim & KNN + PCA
      & KNN pipeline design with PCA, dimensionality analysis, Streamlit dashboard, report Section 3.5 and 4.5
      & 20\% \\
    \bottomrule
  \end{tabularx}
\end{table}

\noindent\textbf{Group contributions} (Sections 1, 2, 5 — marked as group effort):
\begin{itemize}[noitemsep]
  \item Section 1 (Introduction): all members contributed
  \item Section 2 (Data pre-processing): all members contributed; EDA and Notebook 01 jointly authored
  \item Section 5 (Discussion \& Conclusions): all members contributed; final draft consolidated by [lead member]
  \item Notebook 07 (group comparison): all members contributed results CSVs; analysis by all members
\end{itemize}

\subsection*{Individual Declarations}

\noindent I confirm that the contribution statement above accurately represents my individual
contribution to the group work.\\[0.6cm]

\noindent\begin{tabular}{lll}
  \textbf{Name} & \textbf{Signature} & \textbf{Date} \\[0.3cm]
  Gokulkrishna Santoshkumar & \underline{\hspace{5cm}} & \underline{\hspace{3cm}} \\[0.3cm]
  Joel Jolly                 & \underline{\hspace{5cm}} & \underline{\hspace{3cm}} \\[0.3cm]
  Apprajit Vaibhav           & \underline{\hspace{5cm}} & \underline{\hspace{3cm}} \\[0.3cm]
  Vino Mohandas              & \underline{\hspace{5cm}} & \underline{\hspace{3cm}} \\[0.3cm]
  Afzal Hakkim               & \underline{\hspace{5cm}} & \underline{\hspace{3cm}} \\[0.3cm]
\end{tabular}

\subsection*{Group Meetings Logbook}

\begin{table}[H]
  \centering
  \caption{Group meeting logbook}
  \rowcolors{2}{lightgray}{white}
  \begin{tabularx}{\linewidth}{clXXl}
    \toprule
    \textbf{No.} & \textbf{Date} & \textbf{Discussion Topics} & \textbf{Work Done} & \textbf{Attendees} \\
    \midrule
    1 & [date] & Dataset selection, CRISP-DM overview, member model allocation & Dataset agreed; EDA notebook started & All \\
    2 & [date] & EDA findings review, feature engineering decisions & Feature list finalised, Notebook 01 completed & All \\
    3 & [date] & Model baseline results, metric alignment & Common metric framework agreed & All \\
    4 & [date] & Tuned results, DT param\_grid differentiation & All member notebooks complete & All \\
    5 & [date] & Group comparison notebook, dashboard review & Notebook 07 and Streamlit finalised & All \\
    6 & [date] & Report draft review, contribution statements & Report finalised for submission & All \\
    \bottomrule
  \end{tabularx}
\end{table}

\clearpage

%% =============================================================
\section*{References}
\addcontentsline{toc}{section}{References}

\begin{thebibliography}{99}

\bibitem{breiman2001}
  Breiman, L. (2001).
  \textit{Random Forests}.
  Machine Learning, 45(1), 5--32.
  \url{https://doi.org/10.1023/A:1010933404324}

\bibitem{chen2016}
  Chen, T., \& Guestrin, C. (2016).
  XGBoost: A scalable tree boosting system.
  \textit{Proceedings of the 22nd ACM SIGKDD}, 785--794.
  \url{https://doi.org/10.1145/2939672.2939785}

\bibitem{cortes1995}
  Cortes, C., \& Vapnik, V. (1995).
  Support-vector networks.
  \textit{Machine Learning}, 20(3), 273--297.
  \url{https://doi.org/10.1007/BF00994018}

\bibitem{rumelhart1986}
  Rumelhart, D., Hinton, G., \& Williams, R. (1986).
  Learning representations by back-propagating errors.
  \textit{Nature}, 323, 533--536.
  \url{https://doi.org/10.1038/323533a0}

\bibitem{cover1967}
  Cover, T., \& Hart, P. (1967).
  Nearest neighbor pattern classification.
  \textit{IEEE Transactions on Information Theory}, 13(1), 21--27.

\bibitem{pearson1901}
  Pearson, K. (1901).
  On lines and planes of closest fit to systems of points in space.
  \textit{The London, Edinburgh, and Dublin Philosophical Magazine}, 2(11), 559--572.

\bibitem{wirth2000}
  Wirth, R., \& Hipp, J. (2000).
  CRISP-DM: Towards a standard process model for data mining.
  \textit{Proceedings of the 4th International Conference on the Practical Applications
  of Knowledge Discovery and Data Mining}, 29--39.

\bibitem{sklearn}
  Pedregosa, F., et al. (2011).
  Scikit-learn: Machine learning in Python.
  \textit{Journal of Machine Learning Research}, 12, 2825--2830.

\bibitem{iata2023}
  International Air Transport Association. (2023).
  \textit{IATA Annual Review 2023}.
  Geneva: IATA.

\bibitem{matthewsmcc}
  Matthews, B. W. (1975).
  Comparison of the predicted and observed secondary structure of T4 phage lysozyme.
  \textit{Biochimica et Biophysica Acta (BBA)}, 405(2), 442--451.

\end{thebibliography}

\clearpage

%% =============================================================
\appendix
\section*{Appendix A: Python Notebook Structure}
\addcontentsline{toc}{section}{Appendix A: Python Notebook Structure}

\begin{table}[H]
  \centering
  \caption{Python notebook inventory and scope}
  \rowcolors{2}{lightgray}{white}
  \begin{tabularx}{\linewidth}{llX}
    \toprule
    \textbf{Notebook} & \textbf{Type} & \textbf{Content} \\
    \midrule
    \texttt{01\_Data\_Preprocessing\_EDA.ipynb}            & Group    & Data loading, cleaning, feature engineering, EDA visualisations, train/val/test split \\
    \texttt{02\_Member1\_RandomForest\_DecisionTree.ipynb} & Individual & M1: Random Forest baseline + GridSearchCV tuning, M1-DT tuning, evaluation, importance \\
    \texttt{03\_Member2\_XGBoost\_DecisionTree.ipynb}      & Individual & M2: XGBoost baseline + tuning, M2-DT tuning, evaluation, importance \\
    \texttt{04\_Member3\_SVM\_DecisionTree.ipynb}          & Individual & M3: SVM (RBF) baseline + tuning, M3-DT tuning, evaluation, permutation importance \\
    \texttt{05\_Member4\_KNN\_DecisionTree.ipynb}          & Individual & M4: KNN+PCA pipeline + tuning, M4-DT tuning, evaluation, permutation importance \\
    \texttt{06\_Member5\_MLP\_DecisionTree.ipynb}          & Individual & M5: MLP baseline + tuning, M5-DT tuning, loss curve, evaluation \\
    \texttt{07\_Group\_Combined\_Comparison.ipynb}         & Group    & Load all CSVs, group F1/AUC/MCC/Gap/Cost charts, delta charts, final rankings \\
    \bottomrule
  \end{tabularx}
\end{table}

\section*{Appendix B: Data Dictionary}
\addcontentsline{toc}{section}{Appendix B: Data Dictionary}

\begin{longtable}{lll}
  \toprule
  \textbf{Feature} & \textbf{Type} & \textbf{Description} \\
  \midrule
  \endhead
  Route\_ID              & Integer  & Unique route identifier \\
  Origin                 & Categorical & Departure airport code \\
  Destination            & Categorical & Arrival airport code \\
  Distance               & Float & Route distance (km) \\
  Aircraft\_Type         & Categorical & Wide-body, Narrow-body, Regional \\
  Ticket\_Price          & Float & Average ticket price (£) \\
  Load\_Factor           & Float & Average passenger load (0--1) \\
  Competition\_Index     & Float & Market competition score (0--1) \\
  Delay\_Minutes         & Float & Average delay (historical mean) \\
  On\_Time\_Performance  & Float & On-time arrival fraction \\
  Fuel\_Efficiency       & Float & Fuel burn per seat-km \\
  Season                 & Categorical & Off-Peak / Shoulder / Peak \\
  Alliance               & Categorical & Star / SkyTeam / Oneworld / None \\
  Region                 & Categorical & Europe / Asia / North America / etc. \\
  Route\_Category        & Categorical & Domestic / Short-haul / Medium-haul / Long-haul \\
  Demand\_Level          & Categorical & Low / Medium / High \\
  \textbf{Profitable}    & Binary & \textbf{Target: 1 = profitable, 0 = loss-making} \\
  price\_per\_km         & Float & Derived: Ticket\_Price / Distance \\
  delay\_flag            & Binary & Derived: Delay\_Minutes $>$ 30 \\
  is\_weekend            & Binary & Derived: departure on weekend \\
  is\_narrow\_body       & Binary & Derived: Aircraft\_Type = Narrow-body \\
  flight\_month          & Integer & Derived: departure month (1--12) \\
  Season\_Ordinal        & Integer & Derived: Off-Peak=0, Shoulder=1, Peak=2 \\
  Demand\_Level\_Ordinal & Integer & Derived: Low=0, Medium=1, High=2 \\
  \textit{Alliance\_*}   & Binary & OHE columns for alliance membership \\
  \textit{Region\_*}     & Binary & OHE columns for route region \\
  \textit{Route\_Cat\_*} & Binary & OHE columns for route category \\
  \bottomrule
  \caption{Full data dictionary — airline route profitability dataset}
\end{longtable}

\end{document}
